\section{Contesto e motivazione}
Le patologie cardiovascolari rappresentano una delle prime cause di morbilità e mortalità a livello globale; il monitoraggio continuo dell'attività cardiaca tramite segnali ECG è pertanto centrale per la diagnosi precoce e la gestione clinica.

In molti scenari applicativi, l'interpretazione del segnale non può essere affidata esclusivamente a una valutazione umana ex post, poiché il valore informativo del battito risiede anche nella \emph{tempestività} con cui viene rilevata un'anomalia.
Da questo punto di vista, l'analisi automatica dei battiti cardiaci pone sfide tecniche sostanziali:
\begin{itemize}
  \item La necessità di elaborazione \textbf{near-real-time},
  \item La gestione di flussi ad alta velocità,
  \item La robustezza rispetto a fault e perdita di pacchetti,
  \item La capacità di mantenere \emph{consistenza} tra fasi di ingestione, predizione e persistenza.
\end{itemize}

Nel dominio medico i trade-off tra \textbf{precisione} e \textbf{sensibilità} ricoprono un ruolo critico.
I falsi negativi possono portare a omissioni diagnostiche rilevanti, mentre i falsi positivi possono introdurre rumore operativo e aumentare il carico di verifica manuale.

Serve dunque una pipeline che produca classificazioni accurate e sia al contempo \emph{interpretabile} dal punto di vista architetturale, scalabile e adatta a infrastrutture distribuite.
In tale prospettiva, il progetto affronta il problema della classificazione di battiti ECG come un caso di studio per l'integrazione tra \textbf{machine learning classico} e \textbf{data streaming}.

Il contesto clinico rende particolarmente rilevante la distinzione tra \emph{monitoraggio episodico} e \emph{monitoraggio continuo}.
Nel primo caso, l'ECG viene acquisito in ambulatorio o in reparto per intervalli limitati; nel secondo, dispositivi portatili o impiantati generano flussi prolungati che devono essere analizzati mentre arrivano.
Le aritmie non sempre si manifestano durante una registrazione breve: ectopie isolate, pause o tachicardie parossistiche possono comparire a distanza di ore, rendendo insufficiente un'analisi puramente batch su archivi statici.
Un sistema in grado di classificare ogni battito al momento dell'acquisizione consente di segnalare tempestivamente pattern anomali, riducendo il ritardo tra evento fisiologico e possibile intervento clinico.

Dal punto di vista ingegneristico, il problema si colloca all'intersezione tra \textbf{signal processing}, \textbf{machine learning} e \textbf{ingegneria dei sistemi distribuiti}.
La letteratura sulla classificazione automatica del MIT--BIH~\cite{moody2001impact,dechazal2004automatic} dimostra da decenni che modelli supervisionati su feature beat-level possono raggiungere accuratezze elevate; tuttavia, gran parte degli studi si concentra sulla fase di training e valutazione offline, senza integrare ingestion, messaggistica e persistenza in un'unica architettura operativa.
Questa tesi colma, in parte, tale divario proponendo una pipeline che riproduce un ciclo di vita realistico del dato, dalla preparazione storica allo scoring in tempo quasi reale.

\section{Approccio e scelte tecnologiche}\label{sec:approccio-e-scelte-tecnologiche}
In questo lavoro si affronta il problema della classificazione di battiti ECG in un contesto di streaming, privilegiando un'architettura che consenta \textbf{disaccoppiamento} tra ingestion e processing, resilienza ai guasti e scalabilità orizzontale.

La soluzione proposta impiega tecnologie consolidate nello stack Big Data:
\begin{itemize}
  \item \textbf{Scala}: implementazione del sistema e dei componenti della pipeline,
  \item \textbf{Apache Kafka}: broker Pub/Sub per la messaggistica asincrona,
  \item \textbf{Apache Spark}: elaborazione batch (MLlib) e streaming (Structured Streaming),
  \item \textbf{Delta Lake}: sink transazionale per la persistenza dei risultati.
\end{itemize}

Questa combinazione consente di separare in modo netto la responsabilità di acquisizione dei dati dalla responsabilità di inferenza, evitando che il carico di una fase condizioni in maniera diretta l'altra.

La scelta di una soluzione \emph{ibrida batch/streaming} risponde a un'esigenza precisa: i modelli predittivi possono essere addestrati su archivi storici con costi computazionali più elevati, mentre l'inferenza deve operare con latenza contenuta e comportamento deterministico.
Il progetto adotta quindi una pipeline che riproduce un ciclo realistico di produzione dei dati, addestramento e serving, offrendo un contesto coerente per valutare non soltanto la qualità del classificatore, ma anche la solidità dell'infrastruttura sottostante.

Un'alternativa avrebbe potuto consistere in un'architettura monolitica: un unico processo che legge il CSV, addestra il modello e classifica i battiti in memoria senza broker intermedi.
Tale soluzione sarebbe più semplice da implementare, ma non permetterebbe di valutare proprietà fondamentali di un sistema distribuito, come il disaccoppiamento tra velocità di produzione e consumo, la tolleranza ai guasti del consumer e la persistenza transazionale dei risultati.
Analogamente, l'adozione di un framework di serving dedicato (TensorFlow Serving, MLflow Model Serving) avrebbe spostato il focus verso l'inferenza a bassa latenza, sacrificando l'integrazione nativa con Spark per training e streaming analitico.
La scelta dello stack Kafka--Spark--Delta risponde invece a un obiettivo didattico e di ricerca: dimostrare come tecnologie consolidate nel panorama Big Data possano essere composte in una pipeline coerente per dati biomedici strutturati, senza dipendere da servizi cloud proprietari.

\section{Obiettivi}\label{sec:obiettivi}
Gli obiettivi perseguiti sono i seguenti:
\begin{enumerate}
  \item Progettare e implementare una \textbf{pipeline end-to-end} per la classificazione di battiti ECG in near-real-time, rispettando requisiti di tolleranza ai fault, scalabilità e disaccoppiamento tra le componenti.
  \item Sviluppare un processo di \textbf{addestramento batch} per un modello basato su Random Forest, includendo fasi di feature engineering, selezione delle variabili e validazione sperimentale.
  \item Implementare un meccanismo di \textbf{ingestion realistico} tramite un simulatore (\texttt{EcgPatientSimulator}) che converte righe del CSV preprocessato (derivato dal MIT-BIH) in payload JSON e li pubblica su Kafka.
  \item Eseguire \textbf{inferenza in streaming} con Spark Structured Streaming, garantendo atomicità delle scritture su Delta Lake tramite checkpointing e \texttt{foreachBatch}, così da preservare la continuità operativa in presenza di interruzioni.
  \item Valutare la qualità del modello con metriche di classificazione e analizzare \emph{latenze}, \emph{throughput} e comportamento del sistema al variare del volume dei messaggi.
\end{enumerate}

Gli obiettivi sono stati formulati in modo da coprire sia la dimensione \emph{funzionale} (il sistema deve classificare correttamente i battiti) sia quella \emph{non funzionale} (il sistema deve operare con latenza contenuta, recuperare da failure e restare comprensibile come architettura).
La separazione tra training batch e inferenza streaming non è un vincolo imposto dal dataset, bensì una scelta progettuale che riflette la frequenza tipica di aggiornamento dei modelli in ambito clinico: il retraining avviene su scala di giorni o settimane, mentre l'inferenza deve procedere a ogni battito ricevuto.
Analogamente, l'uso di un simulatore anziché di hardware reale consente di concentrare l'attenzione sull'integrazione software, lasciando aperte estensioni future verso dispositivi edge con firmware dedicato.

\section{Contributi}\label{sec:contributi}
Si elencano i contributi principali del lavoro:
\begin{itemize}
  \item \textbf{Realizzazione di una pipeline sperimentale completa} in Scala che integra training batch e inferenza in streaming su dati ECG simulati, mantenendo la separazione tra livelli di acquisizione, elaborazione e persistenza.
  \item \textbf{Approccio pragmatico al feature engineering} su battiti cardiaci (intervalli \texttt{pre-RR}/\texttt{post-RR}, descrittori morfologici del QRS, durate PQ/QT/ST e coefficienti \texttt{qrs\_morph} su due derivazioni) e definizione di uno schema JSON compatto per l'ingestion in Kafka.
  \item Dimostrazione dell'integrazione di Spark Structured Streaming con schema dinamico (\texttt{from\_json}), checkpointing e scritture transazionali su Delta Lake per garantire resilienza, consistenza e ripetibilità delle elaborazioni.
  \item Analisi sperimentale completa della performance del modello (tramite \texttt{MulticlassClassificationEvaluator}) e delle proprietà non funzionali della pipeline, con attenzione a latenze, throughput e stabilità operativa.
  \item Formalizzazione di un caso di studio in cui una stessa architettura può essere riletta sia come sistema di prototipazione accademica sia come base per un futuro deployment in ambito e-Health.
\end{itemize}

Rispetto a implementazioni isolate di classificatori ECG presenti in letteratura, il contributo distintivo risiede nell'\textbf{integrazione end-to-end} piuttosto che nella ricerca del massimo score su un benchmark.
Il lavoro documenta non solo il modello predittivo, ma anche le scelte di configurazione (checkpoint, \texttt{foreachBatch}, schema JSON dinamico), i trade-off architetturali (Lambda ibrida, mono-nodo vs cluster) e le implicazioni cliniche degli errori di classificazione.
Questa impostazione rende la tesi utile come riferimento per chi deve progettare pipeline analitiche in ambito sanitario, dove la correttezza del flusso dati è tanto importante quanto l'accuratezza del classificatore.

\section{Struttura della tesi}\label{sec:struttura-della-tesi}
La trattazione è organizzata come segue:
\begin{itemize}
  \item Il \textbf{Capitolo~\ref{ch:dominio-applicativo-e-dataset}} descrive il dominio applicativo e il dataset utilizzato (MIT-BIH in versione tabellare da Kaggle), includendo considerazioni su feature cliniche, preprocessing del segnale e paradigma Edge Computing.
  \item Il \textbf{Capitolo~\ref{ch:lo-stack-tecnologico}} giustifica le scelte tecnologiche adottate, soffermandosi sui vantaggi di Scala, Spark, Kafka KRaft e Delta Lake nel contesto di pipeline distribuite.
  \item Il \textbf{Capitolo~\ref{ch:progettazione-e-architettura-del-sistema}} illustra l'architettura complessiva e la scomposizione in tre moduli principali: \texttt{EcgModelTrainer} (batch), \texttt{EcgPatientSimulator} (ingestion) e \texttt{EcgStreamingPredictor} (streaming).
  \item Il \textbf{Capitolo~\ref{ch:implementazione-e-dettagli-tecnici}} approfondisce dettagli implementativi, mostrando frammenti di codice e spiegando le scelte di configurazione della ML pipeline e del producer Kafka.
  \item Il \textbf{Capitolo~\ref{ch:analisi-dei-risultati-e-test}} presenta la valutazione sperimentale, confrontando metriche di accuratezza, analisi di falsi positivi/negativi, e misure di latenze e throughput.
  \item Il \textbf{Capitolo~\ref{ch:conclusioni}} espone le valutazioni complessive e le possibili linee di estensione verso modelli deep e integrazione analitica con Delta Lake.
\end{itemize}

I sorgenti del repository (\texttt{EcgModelTrainer}, \texttt{EcgPatientSimulator}, \texttt{EcgStreamingPredictor}) sono richiamati nei capitoli di implementazione e risultati, così da legare teoria e codice effettivamente scritto.

La progressione argomentativa segue un percorso dal generale al particolare: il Capitolo~\ref{ch:dominio-applicativo-e-dataset} fonda il problema nel dominio clinico e nel dataset scelto; il Capitolo~\ref{ch:lo-stack-tecnologico} giustifica le tecnologie; i Capitoli~\ref{ch:progettazione-e-architettura-del-sistema} e~\ref{ch:implementazione-e-dettagli-tecnici} descrivono rispettivamente la progettazione logica e i dettagli implementativi; il Capitolo~\ref{ch:analisi-dei-risultati-e-test} chiude il ciclo con la valutazione sperimentale.
Dove possibile, ogni scelta tecnologica è collegata a un requisito esplicito, in modo che il lettore possa ricostruire il ragionamento progettuale senza dover inferire le motivazioni dalle sole configurazioni riportate nel codice.
